\scanpage{46}
\mspage{3.}
en fait, que les groupes de classes de cycles pour la $\handwritten{\tau}$-équivalence sont de type fini. (De là, on déduit d'ailleurs formellement que $C^m(X)\otimes_{\underline{\mathrm{Q}}}\underline{\mathrm{Q}}_{\unsure{l|1|\ell}}\handwritten{\longrightarrow} H^{2m}(X)(m)$ est injectif, donc que le rang de $C^m(X)$ est majoré par $b_{2m}(X)$). D'autre part, $A$ et $B$ impliquent également que la catégorie des motifs construits par des variétés projectives non singulières est semi-simple, et de façon plus terre à terre, que l'anneau des classes de correspondance algébriques (mod $\handwritten{\tau}$-équivalence et tensorisé par $\underline{\mathrm{Q}}$) de $X$ est une algèbre finie semi-simple sur $\underline{\mathrm{Q}}$, qu'on peut d'ailleurs munir d'une involution et d'une trace satisfaisant les conditions habituelles. On en conclut aussi le résultat suivant, qui peut être considéré dans une certaine mesure comme une généralisation de $A$ (lorsqu'on admet que la variante cohomologique de $A$ est déjà acquise) : Soit $H^{2i}(X)\handwritten{\longrightarrow} H^{2i+2r}(X')$ un homomorphisme défini par une classe de correspondance algébrique, et $C^i(X)\handwritten{\longrightarrow} C^{i+r}(X')$ l'homomorphisme qu'elle définit sur les cycles. Alors un élément de $C^{i+r}(X')$ est dans l'image de $C^i(X)$ si et seulement si son image dans $H^{2i+2r}$ est dans celle de $H^{2i}$. - Notes aussi que le résultat de semi-simplicité signalé plus haut impliquerait le résultat de semi-simplicité analogue pour les opérations du groupe de Galois, dans les conjectures de Tate. Enfin, je me suis à peu près convaincu que $A$ et $B$ impliquent également une théorie raisonnable des variétés abéliennes qui interviennent comme variétés de paramètres pour les familles continues de cycles algébriques, et notamment permettent d'obtenir les relations qu'il faut entre ces ``jacobiennes intermédiaires'' (qui peuvent être conçues comme les ``morceaux de nature algébrique'' des jacobiennes de Weil) et la cohomologie de dimension impaire. Mais il est fort possible que la démonstration de $B$ soit déjà liée à l'introduction de ces variétés abéliennes.

En tous cas, $A$ et $A'$ me semblent, dans un sens, préliminaires à $B$, et il semble raisonnable de commencer par elles. Je te signale la forme équivalente assez suggestive de $A'$ :

\uline{Conjecture} $A''$. Soit $U$ l'ouvert affine (de dimension $n$) $X-Y$, alors $C^i(U)=0$ pour tout $i$ tel que $2i\handwritten{>}n$.
% archive sheet number: 45
